% =====================================================================
%  Figure: the three-state exposure chain under two policies (LEAD).
%  SCHEMATIC, NOT DATA. Transition probabilities are the invented matrices
%  of Section 4.5 (sec:markov-example), rows/cols ordered (N, B, A):
%    P_ref = .95 .04 .01 | .50 .40 .10 | .40 .30 .30
%    P_alg = .92 .06 .02 | .30 .55 .15 | .20 .35 .45
%    P_rh  = .92 .06 .02 | 0   .70 .30 | 0   .45 .55
%  Summaries computed exactly (numpy; they agree with tab:markov-illus),
%  sessions starting in N:
%    stationary  ref (.9049,.0719,.0232)  alg (.7708,.1581,.0711)  rh (0,.6,.4)
%    e_20        ref .0207  alg .0599  rh .1893
%    M_{N,A}     ref 64.0   alg 28.3   rh 15.0;  rh enters {B,A} after 12.5 items
%  Edge width = 4pt x probability (hairline minimum 0.3pt).
%  Requires the house colours defined in preamble.tex.
% =====================================================================
\begin{figure}[t]
\centering
\begin{tikzpicture}[x=1cm,y=1cm,
    st/.style={circle,minimum size=0.8cm,inner sep=0pt,font=\small,line width=0.6pt},
    stN/.style={st,fill=LabNavyLight!55,draw=LabNavyLight,text=LabNavy},
    stB/.style={st,fill=LabNavyMid,draw=LabNavyMid,text=white},
    stA/.style={st,fill=LabAccent,draw=LabAccent,text=white},
    pr/.style={font=\tiny,text=LabInk,fill=white,inner sep=0.7pt,rounded corners=1pt},
    panel/.style={font=\footnotesize\bfseries,text=LabNavy,anchor=south west,inner sep=0pt},
    capt/.style={font=\scriptsize,text=LabGrey,inner sep=1pt},
    stat/.style={font=\scriptsize,text=LabInk,anchor=north west,inner sep=0pt,align=left},
  ]
  % \tedge{from}{to}{prob}{colour}{bend}{label pos}: width proportional to prob
  \newcommand{\tedge}[6]{%
    \pgfmathsetmacro{\lw}{max(0.3,4*#3)}%
    \draw[#4,line width=\lw pt,
          -{Stealth[length=2.6pt 1.15,width=2.6pt 1.4]},shorten >=1pt,shorten <=1pt]
      (#1) to[bend left=#5] node[pr,pos=#6] {$#3$} (#2);}
  % \tloop{node}{prob}{colour}{out}{in}{label offset}
  \newcommand{\tloop}[6]{%
    \pgfmathsetmacro{\lw}{max(0.3,4*#2)}%
    \draw[#3,line width=\lw pt,-{Stealth[length=2.6pt 1.15,width=2.6pt 1.4]}]
      (#1) to[out=#4,in=#5,looseness=5.6] (#1);
    \node[pr] at ($(#1)+(#6)$) {$#2$};}
  % \zedge{from}{to}{bend}: a transition with probability zero (no way back)
  \newcommand{\zedge}[3]{%
    \draw[LabRule,line width=0.5pt,dash pattern=on 1.5pt off 1.2pt,
          -{Stealth[length=2.6pt,width=2.6pt]},shorten >=1pt,shorten <=1pt]
      (#1) to[bend left=#3] node[pr,pos=0.5,text=LabGrey] {$0$} (#2);}
  % \statbar{x0}{piN}{piB}{piA}: stationary distribution as a stacked bar
  \newcommand{\statbar}[4]{%
    \def\bw{4.2}%
    \fill[LabNavyLight!55] (#1,-2.12) rectangle ++({#2*\bw},-0.2);
    \fill[LabNavyMid] ({#1+#2*\bw},-2.12) rectangle ++({#3*\bw},-0.2);
    \fill[LabAccent] ({#1+(#2+#3)*\bw},-2.12) rectangle ++({#4*\bw},-0.2);
    \draw[LabNavy,line width=0.3pt] (#1,-2.12) rectangle ++(\bw,-0.2);
    \node[capt,anchor=east] at (#1-0.03,-2.22) {$\varpi$};}

  % ====================== (a) reference policy ======================
  \begin{scope}[shift={(0,0)}]
    \node[panel] at (-0.4,1.78) {(a)\enspace Reference policy $P_{\mathrm{ref}}$};
    \node[stN] (N) at (0.7,0) {$\mathsf N$};
    \node[stB] (B) at (3.3,1.05) {$\mathsf B$};
    \node[stA] (A) at (3.3,-1.05) {$\mathsf A$};
    \node[capt,anchor=north] at ([yshift=-1pt]N.south) {neutral};
    \node[capt,anchor=south east] at ([xshift=-2pt]B.north west) {borderline};
    \node[capt,anchor=north east] at ([xshift=-2pt]A.south west) {antisemitic};
    \tloop{N}{.95}{LabNavyLight}{203}{157}{-0.86,0.36}
    \tloop{B}{.40}{LabNavyMid}{40}{-10}{0.93,0.34}
    \tloop{A}{.30}{LabAccent}{10}{-40}{0.93,-0.34}
    \tedge{N}{B}{.04}{LabNavyMid}{18}{0.5}
    \tedge{B}{N}{.50}{LabNavyLight}{18}{0.5}
    \tedge{N}{A}{.01}{LabAccent}{-18}{0.5}
    \tedge{A}{N}{.40}{LabNavyLight}{-18}{0.5}
    \tedge{B}{A}{.10}{LabAccent}{35}{0.5}
    \tedge{A}{B}{.30}{LabNavyMid}{35}{0.5}
    \statbar{0.2}{0.9049}{0.0719}{0.0232}
    \node[stat] at (0.2,-2.47) {$\varpi(\mathcal A)=0.023$\\[1pt]
      $e_{20}=0.021$,\enspace $M_{\mathsf N\mathcal A}=64$};
  \end{scope}

  % ====================== (b) algorithmic policy ======================
  \begin{scope}[shift={(5.42,0)}]
    \node[panel] at (-0.4,1.78) {(b)\enspace Algorithmic policy $P_{\mathrm{alg}}$};
    \node[stN] (N) at (0.7,0) {$\mathsf N$};
    \node[stB] (B) at (3.3,1.05) {$\mathsf B$};
    \node[stA] (A) at (3.3,-1.05) {$\mathsf A$};
    \tloop{N}{.92}{LabNavyLight}{203}{157}{-0.86,0.36}
    \tloop{B}{.55}{LabNavyMid}{40}{-10}{0.93,0.34}
    \tloop{A}{.45}{LabAccent}{10}{-40}{0.93,-0.34}
    \tedge{N}{B}{.06}{LabNavyMid}{18}{0.5}
    \tedge{B}{N}{.30}{LabNavyLight}{18}{0.5}
    \tedge{N}{A}{.02}{LabAccent}{-18}{0.5}
    \tedge{A}{N}{.20}{LabNavyLight}{-18}{0.5}
    \tedge{B}{A}{.15}{LabAccent}{35}{0.5}
    \tedge{A}{B}{.35}{LabNavyMid}{35}{0.5}
    \statbar{0.2}{0.7708}{0.1581}{0.0711}
    \node[stat] at (0.2,-2.47) {$\varpi(\mathcal A)=0.071$\enspace
      \textcolor{LabGrey}{($\times3.07$)}\\[1pt]
      $e_{20}=0.060$\enspace\textcolor{LabGrey}{($\times2.89$)},\enspace
      $M_{\mathsf N\mathcal A}=28.3$};
  \end{scope}

  % ====================== (c) rabbit-hole variant ======================
  \begin{scope}[shift={(10.84,0)}]
    \node[panel] at (-0.4,1.78) {(c)\enspace A rabbit hole $P_{\mathrm{rh}}$};
    % the closed class {B, A}
    \begin{scope}[on background layer]
      \fill[LabAccent!9,rounded corners=6pt] (2.66,-1.62) rectangle (4.62,1.62);
      \draw[LabAccent!70,line width=0.5pt,rounded corners=6pt,dash pattern=on 2pt off 1.2pt]
        (2.66,-1.62) rectangle (4.62,1.62);
    \end{scope}
    \node[font=\scriptsize\itshape,text=LabAccent,anchor=north,inner sep=1pt]
      at (3.64,-1.64) {closed class $\{\mathsf B,\mathsf A\}$};
    \node[stN] (N) at (0.7,0) {$\mathsf N$};
    \node[stB] (B) at (3.3,1.05) {$\mathsf B$};
    \node[stA] (A) at (3.3,-1.05) {$\mathsf A$};
    \tloop{N}{.92}{LabNavyLight}{203}{157}{-0.86,0.36}
    \tloop{B}{.70}{LabNavyMid}{40}{-10}{0.93,0.34}
    \tloop{A}{.55}{LabAccent}{10}{-40}{0.93,-0.34}
    \tedge{N}{B}{.06}{LabNavyMid}{18}{0.5}
    \zedge{B}{N}{18}
    \tedge{N}{A}{.02}{LabAccent}{-18}{0.5}
    \zedge{A}{N}{-18}
    \tedge{B}{A}{.30}{LabAccent}{35}{0.5}
    \tedge{A}{B}{.45}{LabNavyMid}{35}{0.5}
    \statbar{0.2}{0}{0.6}{0.4}
    \node[stat] at (0.2,-2.47) {$\varpi(\mathcal A)=0.400$\\[1pt]
      $e_{20}=0.189$,\enspace $M_{\mathsf N\mathcal A}=15.0$};
  \end{scope}

  \node[font=\scriptsize\itshape,text=LabGrey,anchor=north east,inner sep=0pt]
    at (15.46,-3.22) {Schematic: invented transition probabilities, not data};
\end{tikzpicture}
\caption{\textbf{A session as a Markov chain over content states, under two
policies.} Each node is a content state, neutral ($\mathsf N$), borderline ($\mathsf B$) or
antisemitic ($\mathsf A$); each arrow is the probability that an item in one state is
followed by an item in another, and its width is proportional to that
probability (zero transitions are dashed). Arrows are coloured by the state
they lead to. \emph{(a)} A reference policy. \emph{(b)} The algorithmic policy
moves at most $0.2$ of each row's probability, yet long-run antisemitic
exposure $\varpi(\mathcal A)$ triples: the extra mass sent from $\mathsf N$ is small,
but borderline and antisemitic content now lead back to neutral content less
often, so sessions that reach them stay longer (\cref{thm:perturb}).
\emph{(c)} A rabbit-hole variant: borderline and antisemitic content never lead
back to neutral content, so $\{\mathsf B,\mathsf A\}$ is a closed class, a rabbit hole. A
session that starts in $\mathsf N$ enters it after $12.5$ items on average and never
leaves, and long-run exposure is $\varpi(\mathcal A)=0.4$ whatever the entry
probability (\cref{prop:markov}(d)). Bars show the stationary
distribution $\varpi$; $e_{20}$ is the expected share of antisemitic items in a
20-item session that starts in $\mathsf N$, and $M_{\mathsf N\mathcal A}$ the expected number
of further items before the first antisemitic one; ratios in (b) are relative
to (a). \textbf{Schematic:} the transition matrices are the invented ones of
\cref{sec:markov-example}, and the summaries are exact computations from them
(\cref{tab:markov-illus}). They describe no platform.}
\label{fig:markov}
\end{figure}
